\section*{Abstract}

UAVs responding to disasters must recognise hazards and people when individual sensors are blinded by smoke, darkness or rotor noise, yet many aerial fusion models use fixed modality weights. We propose UniAdapFuse, a reliability-aware multimodal fusion framework that gates RGB, thermal and audio tokens with a learned Reliability and Uncertainty Estimator before cross-modal attention, allowing a single model to operate with any subset of sensors. We further introduce the Hazard Prior Gate (HPG), a physics-guided module that injects flame chromaticity, smoke achromaticity and edge cues into a lightweight detector. On a 72{,}997-sample corpus built from six public datasets, its classification branch (AdapFuse-v1) reaches 98.74\% combined macro-F1 for disaster and victim-presence classification. A revised joint-training recipe (phased training, stop-gradient isolation and distillation from AdapFuse-v1) restores classification in the joint model after a 53.71\% collapse, but its detection branch does not learn ($\mathrm{mAP}_{50}=0.0$). On D-Fire, HPG does not improve YOLO26n: it lowers $\mathrm{mAP}_{50}$ by 0.94 points and takes about $1.9\times$ longer to train, while YOLO26s with CLAHE reaches 77.10\% $\mathrm{mAP}_{50}$. Our analysis also exposes the limits of current public data: pseudo-thermal inputs derived from RGB and label-matched audio give fusion no advantage over an RGB-only baseline (98.85\%); the audio clip's source class encodes the fire/clean label; the learned gates do not respond to synthetic corruption; a capture-group-aware split lowers combined macro-F1 to 89.14\%; and source identity determines the victim label. These results establish a benchmark and a set of data requirements, namely synchronised sensors, degradation labels and source-decorrelated splits, for evaluating reliability-aware fusion on UAVs.

\vspace{1cm}
\textbf{Keywords:} Aerial Object Detection, UAV Disaster Response, Multimodal Fusion, Reliability-Aware Fusion, Hazard Prior Gate, Cross-Modal Attention, Multi-Task Learning, Gradient Isolation, Knowledge Distillation, Deep Learning
\pagebreak
